%--------------------------------------------------------------------------------------------------
\chapter{Introduction}
%--------------------------------------------------------------------------------------------------

\section{Background and Motivation}

Flooding is a major natural hazard, responsible for substantial economic losses, social disruption, and loss of life across the globe. With changing climate conditions, flood risks are intensifying in both frequency and magnitude. The Intergovernmental Panel on Climate Change (IPCC) projects a rising likelihood of extreme hydrological events, including both pluvial and fluvial floods, due to higher precipitation variability and rising sea levels.

In Europe, and particularly in Central and Eastern regions such as Slovenia, recent flood events have exposed vulnerabilities in both infrastructure and institutional response systems. The need for robust flood risk models that can inform spatial planning, insurance pricing, and asset management is therefore more pressing than ever. This is especially relevant for stakeholders in the real estate and financial sectors, who must now account for climate-related risks as part of their environmental, social, and governance (ESG) reporting and compliance requirements.

\section{Types of Climate Risk}

In the context of climate-related financial disclosures and scenario analysis, climate risks are commonly categorized into three broad types: \textit{physical}, \textit{transition}, and \textit{liability} risks \cite{moorhouse2024}.

\begin{itemize}
    \item \textbf{Physical risks} arise from the direct consequences of climate variability and long-term climate change. These include acute hazards such as floods, storms, and wildfires, as well as chronic effects such as sea-level rise and shifting temperature regimes. These risks can cause immediate physical damage and long-term degradation of asset value.
    
    \item \textbf{Transition risks} arise from the transition to a low-carbon economy. They include policy changes (e.g., carbon pricing, emissions regulations), technology disruption, evolving market preferences, and reputational risks associated with inaction or misalignment.
    
    \item \textbf{Liability risks} are considered by some \cite{woetzel2020climate} and encapsulate risks arising when those affected by anthoropogenic climate change seek compensation.
\end{itemize}

These categories are increasingly embedded in supervisory and regulatory frameworks, including those developed by the Task Force on Climate-related Financial Disclosures (TCFD) and the Network for Greening the Financial System (NGFS). In the European context, the European Central Bank (ECB) requires significant institutions to assess and disclose their exposure to climate-related and environmental risks, including physical risks, as outlined in its \textit{Guide on climate-related and environmental risks} \cite{ecb2020}. Banks are expected to integrate these risks into their internal governance, risk management, and stress testing frameworks. This thesis focuses exclusively on \textbf{acute physical risks}, with a particular emphasis on riverine flooding and its financial impact on real estate assets.


\section{Climate Risk and the ESG Framework}

As institutional investors and financial regulators increasingly emphasize sustainability, the ESG framework has become a central tool for assessing climate-related risks. The environmental dimension of ESG includes metrics such as energy efficiency, greenhouse gas emissions, biodiversity impact, and—critically—exposure to climate-related hazards.

To operationalize physical risk in quantitative models, it is commonly decomposed into three interacting components:

\begin{itemize}
    \item \textbf{Hazard} – the probability and intensity of a climate-related event (e.g., flood depth for a 100-year return period),
    \item \textbf{Vulnerability} – the expected damage to an asset given the intensity of the hazard,
    \item \textbf{Exposure} – the presence and value of assets located in potentially affected areas.
\end{itemize}

These components form the basis of scenario-based risk models that estimate potential losses under varying climate conditions. However, despite their theoretical appeal, such models often suffer from limited spatial precision, scarce validation against real-world observations, and inconsistent data integration across regions. These issues are particularly acute for flood modeling, where local hydrology and building typologies vary substantially across jurisdictions.

\section{Problem Statement}

While global flood hazard models such as the WRI Aqueduct River Floods provide broad coverage and climate scenario integration, they are often limited in spatial resolution and may not reflect local hydrological conditions. On the other hand, high-resolution local models, such as Slovenia’s \textit{Integralna karta globin} (IKG), offer detailed inundation mapping but are typically fragmented, static (i.e., not climate scenario-aware), and difficult to integrate into broader risk platforms.

Furthermore, flood vulnerability is typically modeled using generalized depth--damage functions derived from empirical or synthetic data. These functions may perform well in aggregate, but their applicability across different building types, geographies, and damage-reporting standards remains uncertain. This creates uncertainty in downstream financial estimates and limits the usefulness of such models for asset-level ESG disclosure.

\section{Research Objectives}

This thesis aims to evaluate the performance of global and local flood risk models with a specific focus on hazard and vulnerability components. The main objectives of the research are:

\begin{enumerate}
    \item To compare the predictive accuracy of flood depth estimates from a global model (WRI Aqueduct) and a high-resolution Slovenian model (IKG) using real flood observations.
    \item To assess how well global depth--damage functions (e.g., from the JRC) reflect reported flood damage across different real estate asset types in Slovenia.
    \item To evaluate whether spatial resolution or proximity to water bodies significantly influence model performance.
    \item To develop and document a data onboarding pipeline that enables integration of heterogeneous local flood maps into the modular \textit{Physrisk} physical risk modeling framework.
\end{enumerate}

\section{Methodological Approach}

The thesis adopts a comparative modeling approach using geospatial and statistical techniques. A validation dataset, compiled from the records of the Slovenian Administration for Civil Protection and Disaster Relief, includes appraised damage and measured flood depths for selected events between 2010 and 2023.

Flood hazard models are evaluated using both regression-based and classification-based metrics to assess depth prediction accuracy. Vulnerability models are tested by comparing observed versus predicted damages using segmented analysis by flood depth and real estate type. All data are processed and harmonized into a common geospatial reference system and stored in Zarr format, enabling scalable access within the \textit{Physrisk} platform.

From a scientific perspective, this work contributes to the empirical validation of global flood hazard and vulnerability models in a local Central European context. From a practical standpoint, it improves the transparency and reusability of flood risk data pipelines, supporting efforts to build ESG-compliant risk disclosure tools.

\section{Structure of the Thesis}

The remainder of the thesis is structured as follows:

\begin{itemize}
    \item \textbf{Chapter 2} presents a review of relevant literature and related work on flood hazard modeling, vulnerability estimation, and ESG-aligned risk frameworks.
    \item \textbf{Chapter 3} describes the study area, datasets, and the methodological framework used for model evaluation and data processing.
    \item \textbf{Chapter 4} details the integration of the local IKG model into the Physrisk framework and presents the architecture of the comparison toolset.
    \item \textbf{Chapter 5} provides experimental results, including quantitative evaluation of model performance, depth--damage function analysis, and financial impact estimates.
    \item \textbf{Chapter 6} discusses the key findings, limitations, and possible future research directions, and concludes the thesis.
\end{itemize}


